% Class summary of 30 EEDI demo students (facts sent to the LLM); slides/data/example_teacher_report.json
\begin{tikzpicture}
\begin{axis}[
  xbar stacked, scale only axis, width=\dimexpr\linewidth-4.0cm\relax, height=\figH, bar width=6pt,
  xmin=0, xmax=30, ytick={1,...,11}, enlarge y limits=0.06,
  yticklabels={Gradient Between Two Co-ordinates, Naming Co-ordinates in 2D, Factorising into a Single Bracket,
               Writing Expressions, Simplifying Expressions (Like Terms), Horizontal and Vertical Lines,
               Substitution into Formula, Solving Linear Inequalities, Midpoint Between Two Co-ordinates,
               Linear Sequences (nth term), Linear Equations},
  xlabel={jumlah siswa (dari 30)}, xmajorgrids,
  legend style={at={(0.5,1.03)}, anchor=south, legend columns=3, font=\tiny, /tikz/every even column/.append style={column sep=4pt}},
  y tick label style={font=\tiny},
]
  \addplot+[fill=gapred, draw=none] coordinates {(0,1) (8,2) (9,3) (12,4) (13,5) (5,6) (12,7) (4,8) (4,9) (6,10) (9,11)};
  \addlegendentry{learning gap (\emph{not\_mastered})}
  \addplot+[fill=okgreen, draw=none] coordinates {(0,1) (8,2) (7,3) (8,4) (8,5) (3,6) (7,7) (2,8) (2,9) (3,10) (4,11)};
  \addlegendentry{\emph{mastered}}
  \addplot+[fill=insuf, draw=none] coordinates {(30,1) (14,2) (14,3) (10,4) (9,5) (22,6) (11,7) (24,8) (24,9) (21,10) (17,11)};
  \addlegendentry{bukti kurang}
\end{axis}
\end{tikzpicture}
